%%%%%%%%%%%%%%%%%%%%%%%%%%%%%%%%%
%
%       EXERCÍCIO
%
%%%%%%%%%%%%%%%%%%%%%%%%%%%%%%%%%

\ifdefstring{\atividade}{prova}{%
    \ifdefstring{\modo}{objetivo}{%
        \renewcommand{\valorquestao}{\ValorQObj\ ponto}
    }{%
        \renewcommand{\valorquestao}{\ValorQDisc\ pontos}
    }%
}%

\begin{exercicioBanco}[\valorquestao]
Assinale verdadeira (V) ou falso (F) a cada uma das seguintes sentenças.
%
\begin{enumerate}[\(\bullet\)]
    \item A função \(f(x) = 3,5 - 0,4x\) é decrescente.
    %
    \item O zero de \(f(x) = \dfrac{1}{2}x - 2\) é igual a \(2\).
    %
    \item A função \(f(x) = x - 1\) é negativa no intervalo \(]-\infty,1]\).
    %
    \item O coeficiente angular de \(f(x) = 1 - x\) é igual a \(1\).
\end{enumerate}

% Define as alternativas
\newcommand{\alternativas}{%
    \begin{center}
        \begin{tabularx}{\textwidth}{XXXXX}
            (a) (V, F, V, V). &
            (b) (V, V, F, F). &
            (c) (V, F, V, F). &
            (d) (F, V, F, F). &
            (e) (F, V, F, V).
        \end{tabularx}
    \end{center}
}

% Define a resposta correta
\newcommand{\resposta}{C}

% Lógica condicional para exibição
\ifdefstring{\atividade}{lista}{%
    \alternativas
    \vspace{0.5em}
    
    \noindent\textbf{Resposta:} letra \textbf{\resposta}.
}{%
    \ifdefstring{\modo}{objetivo}{%
        \alternativas
    }{%
        % modo = discursiva → não mostra alternativas
    }
}
\end{exercicioBanco}

